\section{Firmware-AI Integration Workflow}
\label{sec:integration}

Once the neural network code is ready (Workflow 2) and the firmware scaffolding is in place (Workflow 1), Workflow 3 merges the two outputs into a single deployable project. The workflow handles path collection, firmware layout detection, AI file transfer, \texttt{main.c} patching, and a final consistency check, all without manual intervention.

The implementation resides in \texttt{workflow3\_integration.py} and runs five LangGraph nodes in sequence: \texttt{collect\_integration\_info}, \texttt{scan\_ai\_files}, \texttt{copy\_ai\_files}, \texttt{modify\_main\_c}, and \texttt{verify\_integration}.

\subsection{Project Path Collection and Idempotency}

The \texttt{collect\_integration\_info()} node determines two key paths: the firmware project directory and the folder containing the AI-generated code. When the workflow is triggered directly by the user (not as part of a sequential WF1→WF2→WF3 chain), these paths are extracted from a natural language message using a Pydantic schema (\texttt{IntegrationInfoExtraction}) with structured LLM output.

If the paths are already present in the graph state from a prior step, the node skips collection entirely and goes directly to validation. This avoids unnecessary LLM calls when the workflow is invoked automatically.

When the paths cannot be inferred from the message, the system reads the agent's persistent memory to suggest the last known project directory:

\begin{lstlisting}[language=Python, caption={Persistent context suggestion during path collection}, label={lst:integration_suggest}]
last_fw = state.persistent_context.get("last_project_path", "Nessuno")
dynamic_prompt = {
    "instruction": "...",
    "suggestion": f"Last project was: {last_fw}. Use the same or a new path?"
}
resume_value = interrupt(dynamic_prompt)
\end{lstlisting}

If the user's reply contains a word like ``yes'', ``previous'', or ``profile'', the path from the persistent context is applied directly, with no further LLM extraction.

\subsection{Filesystem Validation and Layout Detection}

Before any file operation begins, the internal helper \texttt{\_validate\_and\_detect\_structure()} expands relative and home-directory references (e.g., \texttt{\textasciitilde/} or \texttt{./}) using \texttt{os.path.expanduser()} and checks that both paths exist on disk. Path expansion happens before any scan so that downstream functions always work with absolute paths.

The node then detects which of the two common STM32 project layouts is in use. The standard layout places source files under \texttt{Src/} and headers under \texttt{Inc/}, while the STM32Cube layout uses \texttt{Core/Src/} and \texttt{Core/Inc/}. The correct pair of directories is selected automatically. If neither pattern matches, the workflow raises an error before touching any file.

The node also handles the case where the project root contains a single subdirectory: in that case, the subdirectory is treated as the real project root, which is common when a CubeMX project is zipped and extracted into a wrapper folder.

\subsection{AI File Scanning}

The \texttt{scan\_ai\_files()} node lists the contents of the AI code directory and splits the files into two groups: \texttt{.c} sources (stored in \texttt{state.ai\_src\_files}) and \texttt{.h} headers (stored in \texttt{state.ai\_header\_files}). If both groups are empty, the workflow raises an error rather than producing a silent, broken integration.

\subsection{File Placement and X-CUBE-AI Runtime Headers}
\label{sec:copy_ai_files}

The \texttt{copy\_ai\_files()} node copies the scanned files into the firmware directories identified earlier: \texttt{.c} files go to \texttt{firmware\_src\_dir} and \texttt{.h} files go to \texttt{firmware\_inc\_dir}.

Beyond the network files, the node also handles a less obvious requirement: the X-CUBE-AI runtime headers. When STEdgeAI generates the network code, it also produces a subdirectory at \texttt{st\_ai\_ws/inspector\_network/workspace/include} containing the STMicroelectronics AI platform headers (such as \texttt{ai\_platform.h}). These files are needed at compile time, but are not included in a standard CubeMX project unless the X-CUBE-AI component has been explicitly added through the \texttt{.ioc} configuration.

The node checks for this directory and, if it exists, copies all \texttt{.h} files into \texttt{Middlewares/ST/AI/Inc/}, creating the folder if necessary:

\begin{lstlisting}[language=Python, caption={X-CUBE-AI runtime header copy}, label={lst:xcubeai_headers}]
ws_include_dir = os.path.join(output_root, "..", "st_ai_ws",
    "inspector_network", "workspace", "include")
middlewares_ai_inc = os.path.join(proj_root, "Middlewares", "ST", "AI", "Inc")

if os.path.exists(ws_include_dir):
    os.makedirs(middlewares_ai_inc, exist_ok=True)
    for header in os.listdir(ws_include_dir):
        if header.endswith('.h'):
            shutil.copy2(src, dst)
else:
    logger.warning("Runtime headers not found. Add X-CUBE-AI via CubeMX .ioc.")
\end{lstlisting}

If the directory is absent, a warning is logged and the user is advised to add the component through CubeMX. This check prevents the common compile-time error caused by a missing \texttt{ai\_platform.h}.

\subsection{Automated \texttt{main.c} Modification}
\label{sec:main_c_injection}

The \texttt{modify\_main\_c()} node is the most involved step. It reads the existing \texttt{main.c} and injects AI initialization and inference code into the STM32CubeMX user-code comment blocks (\texttt{USER CODE BEGIN}). Since CubeMX regenerates only its own auto-generated regions, code placed inside these blocks is preserved across future regenerations.

Before writing anything, the node creates a timestamped backup:

\begin{lstlisting}[language=Python, caption={Backup before patching}, label={lst:main_backup}]
backup_path = f"{state.main_c_path}.backup_{state.timestamp}"
shutil.copy2(state.main_c_path, backup_path)
\end{lstlisting}

Four injections are then applied in sequence, each guarded by a regex check to skip the block if the corresponding code is already present:

\begin{enumerate}[noitemsep, topsep=4pt]
    \item \textbf{USER CODE BEGIN Includes}: adds \texttt{\#include "network.h"} and \texttt{"network\_data.h"}.
    \item \textbf{USER CODE BEGIN PV}: declares the AI handle, the activations buffer, and the input/output float arrays (\texttt{in\_data}, \texttt{out\_data}).
    \item \textbf{USER CODE BEGIN 2}: calls \texttt{ai\_network\_create()}, \texttt{ai\_network\_init()}, and binds the buffer pointers.
    \item \textbf{USER CODE BEGIN WHILE}: calls \texttt{ai\_network\_run()} in the main loop, with commented examples for sensor data input and output classification.
\end{enumerate}

The injected loop code is shown below:

\begin{lstlisting}[language=C, caption={Injected inference code in the main loop (USER CODE BEGIN WHILE)}, label={lst:main_while}]
/* Fill in_data with sensor data (ADC, I2C, SPI, etc.)
   HAL_ADC_Start(&hadc1);
   HAL_ADC_PollForConversion(&hadc1, HAL_MAX_DELAY);
   in_data[0] = (ai_float)HAL_ADC_GetValue(&hadc1) / 4095.0f;
*/

if (ai_network_run(network, &ai_input[0], &ai_output[0]) != 1) {
    /* Inference failed - blink error LED */
    HAL_GPIO_WritePin(LED_ERROR_GPIO_Port, LED_ERROR_Pin, GPIO_PIN_SET);
    HAL_Delay(100);
    HAL_GPIO_WritePin(LED_ERROR_GPIO_Port, LED_ERROR_Pin, GPIO_PIN_RESET);
}

/* Process results (classification example):
   int predicted_class = 0;
   ai_float max_prob = out_data[0];
   for (int i = 1; i < AI_NETWORK_OUT_1_SIZE; i++) {
       if (out_data[i] > max_prob) {
           max_prob = out_data[i];
           predicted_class = i;
       }
   }
*/
\end{lstlisting}

The network name comes from \texttt{state.network\_name}, so the injection works across different model names. If \texttt{modify\_main} is set to \texttt{False} in the state, the step is skipped cleanly.

\subsection{Verification and Persistent Context Update}

The \texttt{verify\_integration()} node checks that every copied file is present at its destination and that the \texttt{main.c} patch was recorded as successful. If all checks pass, the workflow saves the used paths to Redis (see Section~\ref{sec:redis_architecture}):

\begin{lstlisting}[language=Python, caption={Persistent context update on success}, label={lst:integration_persist}]
state.persistent_context["last_project_path"] = state.firmware_project_dir
state.persistent_context["last_ai_code_dir"] = state.ai_code_dir
\end{lstlisting}

These values are then available in future sessions, so the user will be prompted with the same paths as a default suggestion the next time the workflow runs.

If verification fails, the response includes a step-by-step breakdown showing which operations succeeded and which did not, to help with manual recovery.

\subsection{Build Verification}

As a final check, the workflow compiles the project with the configured toolchain (typically \texttt{arm-none-eabi-gcc}). A clean build confirms that the AI files are in the right place, the runtime headers are visible to the compiler, and the injected code in \texttt{main.c} has no syntax errors. If the build fails, the compiler log is returned directly to the user for inspection.
